\section{Recursos}
\label{sec:recursos}

\vspace{-0.15cm}

\subsection{Recursos humanos}

\textbf{Realización del proyecto y documentación:}
\begin{itemize}
    \item \textbf{Nombre:} \AutorTFG
    \item \textbf{Email:} \texttt{\EmailTFG}
    \item \textbf{Titulación:} \TitulacionTFG. \\\MencionTFG.
\end{itemize}

\textbf{Directores del proyecto:}
\begin{itemize}
    \item \DirectorUno. Profesor Permanente Laboral del Departamento de Ciencia de la Computación e Inteligencia Artificial de la Universidad de Córdoba, Escuela Politécnica Superior de Córdoba. Miembro investigador del grupo AYRNA.
    \item \DirectorDos. Profesor Sustituto del Departamento de Ciencia de la Computación e Inteligencia Artificial de la Universidad de Córdoba, Escuela Politécnica Superior de Córdoba. Miembro investigador del grupo AYRNA.
\end{itemize}

\vspace{-0.15cm}

\subsection{Recursos software}
\begin{itemize}
    \item \textbf{Sistema operativo:} \texttt{Ubuntu 22.04 LTS 64 bits}.
    \item \textbf{Entorno de desarrollo (\texttt{IDE}):} \texttt{Visual Studio Code}.
    \item \textbf{Entorno de desarrollo web local:} Servidor web integrado de \texttt{Spring Boot} (para el \textit{backend}) y \texttt{Node.js} (para el \textit{frontend}).
    \item \textbf{Plataforma de despliegue y evaluación:} \texttt{Docker} y \texttt{Docker Compose}.
    \item \textbf{Gestor de base de datos:} \texttt{MongoDB} y \texttt{MongoDB Compass}.
    \item \textbf{Metodología de pruebas:} Pruebas funcionales y de 
    usabilidad manuales en diferentes navegadores web.
\end{itemize}

%\begin{samepage}
\vspace{-0.15cm}

\subsection{Recursos hardware}
Se utilizará un ordenador con las siguientes características:
\begin{itemize}
    \item \textbf{Procesador:} Intel\textregistered{} Core\texttrademark{} i5-1035G1 CPU @ 1.00GHz (8 CPUs), 1.20 GHz.
    \item \textbf{Memoria RAM:} 8GB.
    \item \textbf{Gráfica:} Intel\textregistered{} UHD Graphics
\end{itemize}


%\end{samepage}